% GENERATED FILE. Edit canonical lessons, then run npm run generate:books.
% canonical-source-hash: fnv1a64:ae8864e837170f4f
% canonical-lessons: IT-C518-intonato, IT-C518-stonare, IT-C518-ritornello, IT-C518-strofa, IT-C518-canticchiare, IT-R518-ripasso-511-518

\chapter{Intonati e stonati}
\label{ch:it-intonati-e-stonati}
\hlchaptermodality{518}

\begin{chapteropening}
\textbf{By the end of this chapter:} \emph{I can talk about singing in tune and out of tune.}

Complete the last lesson of chapter 518: I can talk about singing in tune and out of tune.
\end{chapteropening}

\section[intonato]{intonato — in tune}

\label{lesson:IT-C518-intonato}



\begin{quote}
Before the new one: say the Italian for to tune, then the Italian for to compose.
\end{quote}

\subsection*{You'll want to know: intonato}
\textbf{intonato} — \textquotedblleft{}in tune\textquotedblright{}.

\begin{quote}
Hai una voce molto intonata.
\end{quote}

\begin{quote}
\emph{You sing very much in tune.}
\end{quote}

Also matching: \emph{una borsa intonata alle scarpe}.

\begin{cousinweb}
\emph{intonato} comes from \emph{intonare} \textquotedblleft{}to tune, to strike up\textquotedblright{}, built on \emph{tono} \textquotedblleft{}tone\textquotedblright{}, from Latin \emph{tonus}, from Greek \emph{tónos} \textquotedblleft{}a stretching, a pitch\textquotedblright{} (EtymDB).
\end{cousinweb}

\begin{grammarlens}[title={What you've built}]
One more word for telling, explaining and describing at B1.
\end{grammarlens}

\subsection*{Your turn}
\begin{itemize}
\raggedright
  \item \emph{Say it:} \emph{intonato}
  \item \emph{Say it:} \emph{intonato}, once more
  \item \emph{Say it:} a sentence of your own with \emph{intonato}
  \item \emph{Recall:} say the Italian for to tune, then the Italian for to compose, then say \emph{intonato} again
\end{itemize}

\begin{tcolorbox}[breakable,colback=teal!4,colframe=teal!35!black,title={Before you move on}]
What does intonato mean? (\textquotedblleft{}In tune\textquotedblright{}.) Say it once more.
\end{tcolorbox}
\section[stonare]{stonare — to sing or play out of tune}

\label{lesson:IT-C518-stonare}



\begin{quote}
Before the new one: say the Italian for to compose, then the Italian for in tune.
\end{quote}

\subsection*{You'll want to know: stonare}
\textbf{stonare} — \textquotedblleft{}to sing or play out of tune\textquotedblright{}.

\begin{quote}
Quando canto stono sempre.
\end{quote}

\begin{quote}
\emph{When I sing I am always out of tune.}
\end{quote}

\emph{Stonato} — out of tune; also out of place.

\begin{cousinweb}
\emph{stonare} is built from \emph{s-} and \emph{tono} \textquotedblleft{}tone\textquotedblright{}, from Latin \emph{tonus}, from Greek \emph{tónos} \textquotedblleft{}a pitch\textquotedblright{} (EtymDB): to go off the note.
\end{cousinweb}

\begin{grammarlens}[title={What you've built}]
One more word for telling, explaining and describing at B1.
\end{grammarlens}

\subsection*{Your turn}
\begin{itemize}
\raggedright
  \item \emph{Say it:} \emph{stonare}
  \item \emph{Say it:} \emph{stonare}, once more
  \item \emph{Say it:} a sentence of your own with \emph{stonare}
  \item \emph{Recall:} say the Italian for to compose, then the Italian for in tune, then say \emph{stonare} again
\end{itemize}

\begin{tcolorbox}[breakable,colback=teal!4,colframe=teal!35!black,title={Before you move on}]
What does stonare mean? (\textquotedblleft{}To sing or play out of tune\textquotedblright{}.) Say it once more.
\end{tcolorbox}
\section[il ritornello]{il ritornello — chorus, refrain}

\label{lesson:IT-C518-ritornello}



\begin{quote}
Before the new one: say the Italian for in tune, then the Italian for to sing or play out of tune.
\end{quote}

\subsection*{You'll want to know: il ritornello}
\textbf{il ritornello} — \textquotedblleft{}chorus, refrain\textquotedblright{}.

\begin{quote}
Tutti cantano il ritornello insieme.
\end{quote}

\begin{quote}
\emph{Everyone sings the chorus together.}
\end{quote}

Also a phrase someone repeats: \emph{sempre lo stesso ritornello}.

\begin{cousinweb}
\emph{ritornello} is built on \emph{ritorno} \textquotedblleft{}return\textquotedblright{}, from \emph{ritornare}, from Latin \emph{tornāre} \textquotedblleft{}to turn on a lathe\textquotedblright{} (Lewis \& Short): the part that comes back.
\end{cousinweb}

\begin{grammarlens}[title={What you've built}]
One more word for telling, explaining and describing at B1.
\end{grammarlens}

\subsection*{Your turn}
\begin{itemize}
\raggedright
  \item \emph{Say it:} \emph{il ritornello}
  \item \emph{Say it:} \emph{il ritornello}, once more
  \item \emph{Say it:} a sentence of your own with \emph{il ritornello}
  \item \emph{Recall:} say the Italian for in tune, then the Italian for to sing or play out of tune, then say \emph{il ritornello} again
\end{itemize}

\begin{tcolorbox}[breakable,colback=teal!4,colframe=teal!35!black,title={Before you move on}]
What does il ritornello mean? (\textquotedblleft{}Chorus, refrain\textquotedblright{}.) Say it once more.
\end{tcolorbox}
\section[la strofa]{la strofa — verse, stanza}

\label{lesson:IT-C518-strofa}



\begin{quote}
Before the new one: say the Italian for to sing or play out of tune, then the Italian for chorus.
\end{quote}

\subsection*{You'll want to know: la strofa}
\textbf{la strofa} — \textquotedblleft{}verse, stanza\textquotedblright{}.

\begin{quote}
Non ricordo la seconda strofa della canzone.
\end{quote}

\begin{quote}
\emph{I do not remember the second verse of the song.}
\end{quote}

Also \emph{la strofe}.

\begin{cousinweb}
\emph{strofa} is a learned borrowing, through Latin \emph{stropha}, from Greek \emph{stroph\'{\={e}}} \textquotedblleft{}a turning\textquotedblright{} (EtymDB, under English \emph{strophe}), the turn of the chorus as it danced.
\end{cousinweb}

\begin{grammarlens}[title={What you've built}]
One more word for telling, explaining and describing at B1.
\end{grammarlens}

\subsection*{Your turn}
\begin{itemize}
\raggedright
  \item \emph{Say it:} \emph{la strofa}
  \item \emph{Say it:} \emph{la strofa}, once more
  \item \emph{Say it:} a sentence of your own with \emph{la strofa}
  \item \emph{Recall:} say the Italian for to sing or play out of tune, then the Italian for chorus, then say \emph{la strofa} again
\end{itemize}

\begin{tcolorbox}[breakable,colback=teal!4,colframe=teal!35!black,title={Before you move on}]
What does la strofa mean? (\textquotedblleft{}Verse, stanza\textquotedblright{}.) Say it once more.
\end{tcolorbox}
\section[canticchiare]{canticchiare — to hum, to sing softly}

\label{lesson:IT-C518-canticchiare}



\begin{quote}
Before the new one: say the Italian for chorus, then the Italian for verse.
\end{quote}

\subsection*{You'll want to know: canticchiare}
\textbf{canticchiare} — \textquotedblleft{}to hum, to sing softly\textquotedblright{}.

\begin{quote}
Canticchio mentre cucino.
\end{quote}

\begin{quote}
\emph{I hum while I cook.}
\end{quote}

Built on \emph{cantare} with \emph{-icchiare}, for a light, repeated action.

\begin{cousinweb}
\emph{canticchiare} is built on \emph{cantare} \textquotedblleft{}to sing\textquotedblright{} (EtymDB), inherited from Latin \emph{cantāre}.
\end{cousinweb}

\begin{grammarlens}[title={What you've built}]
One more word for telling, explaining and describing at B1.
\end{grammarlens}

\subsection*{Your turn}
\begin{itemize}
\raggedright
  \item \emph{Say it:} \emph{canticchiare}
  \item \emph{Say it:} \emph{canticchiare}, once more
  \item \emph{Say it:} a sentence of your own with \emph{canticchiare}
  \item \emph{Recall:} say the Italian for chorus, then the Italian for verse, then say \emph{canticchiare} again
\end{itemize}

\begin{tcolorbox}[breakable,colback=teal!4,colframe=teal!35!black,title={Before you move on}]
What does canticchiare mean? (\textquotedblleft{}To hum, to sing softly\textquotedblright{}.) Say it once more.
\end{tcolorbox}
\section[(dialogue)]{Ripasso dei capitoli 511–518}

\label{lesson:IT-R518-ripasso-511-518}



\begin{quote}
Say the Italian for to score and draw.
\end{quote}

\subsection*{Your turn}
Say these aloud:

\begin{itemize}
\raggedright
  \item to score, draw, to draw, league table, to knock out
  \item tournament, semi-final, to qualify, medal, podium
  \item to compete, stopwatch, marathon, athlete, to pant
  \item dive, swimming cap, sail, to dive under, diving board
  \item to skate, boot, slope, skier, chairlift
  \item cello, flute, trumpet, accordion, to strum
  \item orchestra, symphony, sheet music, to tune, to compose
  \item in tune, to sing or play out of tune, chorus, verse, to hum
\end{itemize}

\begin{tcolorbox}[breakable,colback=teal!4,colframe=teal!35!black,title={Before you move on}]
To score? (\textbf{segnare}.) Draw? (\textbf{il pareggio}.) To draw? (\textbf{pareggiare}.) League table? (\textbf{la classifica}.) To knock out? (\textbf{eliminare}.) Tournament? (\textbf{il torneo}.) Semi-final? (\textbf{la semifinale}.) To qualify? (\textbf{qualificarsi}.) Medal? (\textbf{la medaglia}.) Podium? (\textbf{il podio}.) To compete? (\textbf{gareggiare}.) Stopwatch? (\textbf{il cronometro}.) Marathon? (\textbf{la maratona}.) Athlete? (\textbf{l'atleta}.) To pant? (\textbf{ansimare}.) Dive? (\textbf{il tuffo}.) Swimming cap? (\textbf{la cuffia}.) Sail? (\textbf{la vela}.) To dive under? (\textbf{immergersi}.) Diving board? (\textbf{il trampolino}.) To skate? (\textbf{pattinare}.) Boot? (\textbf{lo scarpone}.) Slope? (\textbf{la pista}.) Skier? (\textbf{lo sciatore}.) Chairlift? (\textbf{la seggiovia}.) Cello? (\textbf{il violoncello}.) Flute? (\textbf{il flauto}.) Trumpet? (\textbf{la tromba}.) Accordion? (\textbf{la fisarmonica}.) To strum? (\textbf{strimpellare}.) Orchestra? (\textbf{l'orchestra}.) Symphony? (\textbf{la sinfonia}.) Sheet music? (\textbf{lo spartito}.) To tune? (\textbf{accordare}.) To compose? (\textbf{comporre}.) In tune? (\textbf{intonato}.) To sing or play out of tune? (\textbf{stonare}.) Chorus? (\textbf{il ritornello}.) Verse? (\textbf{la strofa}.) To hum? (\textbf{canticchiare}.)
\end{tcolorbox}
